\chapter{Mômen động lượng, spin và thống kê}

\section{Từ sự quay thông thường đến spin nội tại}

Mômen động lượng là chủ đề hình thức, nhưng spin mới là đích thực sự. Spin là một trong những tính chất đơn giản nhất mà cũng ít trực giác nhất của các hạt cơ bản, và cấu trúc của nó suy ra từ toán học của mômen động lượng.

Mômen động lượng là một vectơ. Hướng của nó nằm dọc theo trục quay. Có hai hướng khả dĩ dọc theo trục ấy, nên cần một quy ước để chọn giữa chúng: cuộn các ngón tay phải theo chiều quay, và ngón cái sẽ chỉ theo hướng của vectơ mômen động lượng. Đây là một định nghĩa, không phải định lý.

Độ lớn phụ thuộc vào lượng khối lượng hiện diện, cách khối lượng ấy phân bố, kích thước của hệ và vận tốc góc của nó. Sự phân bố khối lượng và kích thước được gộp vào mômen quán tính. Vì vậy một vật lớn và một vật nhỏ có thể mang cùng một mômen động lượng trong khi quay với tốc độ rất khác nhau và có năng lượng rất khác nhau.

Một vật hợp thành thông thường có thể mang hai loại mômen động lượng khác nhau về mặt khái niệm. Một chiếc cốc chuyển động theo đường tròn quanh một trục bên ngoài có mômen động lượng vì khối tâm của nó chuyển động quanh trục đó, ngay cả khi chiếc cốc không quay quanh bất kỳ trục nào đi qua chính nó. Đây là mômen động lượng quỹ đạo. Spin có nghĩa là sự quay quanh một trục bên trong.

Định nghĩa sáng sủa nhất dùng hệ quy chiếu khối tâm. Xét hệ quy chiếu mà trong đó khối tâm đứng yên, nên tổng động lượng thông thường triệt tiêu. Mọi mômen động lượng còn lại trong hệ quy chiếu đó đều là mômen động lượng spin. Vật không nhất thiết phải đứng yên về mặt nội tại; chỉ có khối tâm của nó là đứng yên.

Đối với một quả bóng rổ, sự phân biệt này không gây khó khăn gì. Một quả bóng rổ đứng yên và chính quả bóng ấy được cho quay rõ ràng vẫn được nhận ra là cùng một vật ở hai trạng thái khác nhau. Cơ học lượng tử làm bức tranh này phức tạp hơn vì mômen động lượng là rời rạc. Không có dãy liên tục nào của các mômen động lượng được phép nối một trạng thái có mômen động lượng bằng 0 với trạng thái quay đầu tiên. Vì vậy người ta có thể hỏi liệu trạng thái quay đầu tiên có phải là cùng một vật nhưng đang chuyển động, hay hoàn toàn là một vật khác.

Cách diễn đạt phần nào mang tính quy ước, nhưng câu hỏi vật lý thì chính xác: cần bao nhiêu năng lượng để đạt tới trạng thái quay đầu tiên khác không? Nếu khoảng cách mức nhỏ, vật bị kích thích vẫn dễ dàng được nhận ra. Nếu khoảng cách ấy cực lớn, sự kích thích có thể làm thay đổi vật một cách căn bản. Việc cố làm cho một nguyên tử có kích thích quay quá lớn có thể làm nó văng mất các electron và phá hủy nguyên tử, thay vì để chính nguyên tử ấy quay một cách êm ả.

Bây giờ hãy xét một electron, tạm thời gác sang một bên spin nội tại đã biết của nó. Có thể làm cho một electron quay quanh một trục bên trong không? Một điểm toán học không có cấu trúc bên trong nhìn thấy được để quay, nhưng thực nghiệm không xác lập kích thước của electron bằng cách cung cấp một bề mặt cứng nhỏ mà ta có thể theo dõi khi nó quay. Nếu electron có kích thước hữu hạn nhưng cực kỳ nhỏ, thì kích thích quay bổ sung đầu tiên của nó đơn giản có thể đòi hỏi năng lượng quá lớn mới đạt tới được. Vì năng lượng tăng lên khi kích thước giảm xuống, năng lượng cần thiết có thể lớn đến mức thiên văn, thậm chí có lẽ còn cỡ năng lượng Planck. Ước lượng này cố ý chỉ có tính tạm thời.\footnote{Một đoạn ngắn ngay sau ước lượng này không đủ dễ hiểu để làm cơ sở cho việc tái dựng thêm nội dung thực chất.}

Vì thế, đối với một hạt cơ bản, độ lớn của spin được xem là một đặc trưng cố định. Electron không đi kèm một dải các độ lớn spin nội tại tăng dần liên tiếp. Nếu một kích thích về mọi phương diện khác đều giống electron nhưng mang tổng spin khác, nó sẽ được phân loại là một hạt khác.

Độ lớn cố định không ấn định hướng. Hãy tưởng tượng một hạt gọi là \emph{spintron}, được nhận diện một phần bởi độ lớn mômen động lượng của nó. Nếu spin của nó chỉ dọc theo một trục, thì tính bất biến quay đòi hỏi rằng cùng loại hạt ấy cũng có thể có spin chỉ dọc theo một trục khác. Không có hướng không gian nào được ưu tiên một cách nội tại. Độ lớn spin vẫn cố định trong khi định hướng của nó có thể thay đổi. Vì vậy bài toán có hai thành phần dường như xung đột nhau: một độ lớn bị lượng tử hóa và những hướng có thể quay tự do. Đại số mômen động lượng sẽ chỉ ra cách chúng cùng tồn tại.

\begin{classroomqa}
\lecturetimestamp{00:12:25}
\audiencequestion{Vì sao cần năng lượng để làm một hạt nhỏ quay?}
\lecturerresponse{Đối với một vật có khối lượng \(M\) và bán kính đặc trưng \(R\), mômen quán tính có bậc cỡ
\begin{equation}
I\sim cMR^2,
\end{equation}
trong đó hệ số số học \(c\) phụ thuộc vào hình dạng và vào cách khối lượng được phân bố. Nếu phần vật chất gần phía ngoài chuyển động với tốc độ \(v\), thì động lượng thông thường của nó có bậc cỡ \(Mv\), và mômen động lượng bằng động lượng nhân khoảng cách:
\begin{equation}
L\sim MvR.
\end{equation}
Tương đương, \(L=I\omega\) với \(v\sim\omega R\). Động năng quay là
\begin{equation}
E_{\mathrm{rot}}=\frac{L^2}{2I}.
\end{equation}
Khi cố định \(M\) và cố định \(L\),
\begin{equation}
E_{\mathrm{rot}}\sim\frac{L^2}{2cMR^2}\propto\frac{1}{R^2}.
\end{equation}
Vật càng nhỏ thì một mômen động lượng cho trước càng tốn nhiều năng lượng. Sự phụ thuộc theo nghịch đảo bình phương đó là cổ điển. Sự kiện bổ sung của cơ học lượng tử là mômen động lượng thay đổi theo các bước rời rạc, nên bước được phép đầu tiên không thể làm nhỏ tùy ý.}
\end{classroomqa}

Phần toán học sắp dùng thì trừu tượng và lúc đầu khá ít trực giác. Điều làm nó đáng quan tâm là các sự kiện có thể kiểm nghiệm bằng thực nghiệm lại xuất hiện từ một tập nhỏ các quy tắc đại số.

\section{Mômen động lượng quỹ đạo}

Hãy bắt đầu với một hạt chuyển động trên quỹ đạo tròn quanh một gốc tọa độ. Mômen động lượng của nó có độ lớn bằng động lượng nhân với khoảng cách vuông góc tính từ gốc. Trong một mặt phẳng, ta có thể gọi một chiều quay là dương và chiều ngược lại là âm. Dấu ấy ghi lại hướng của vectơ mômen động lượng so với trục đã chọn.

Đây vẫn chưa phải là spin. Hạt có mômen động lượng vì nó có động lượng thông thường. Bây giờ đặt hai hạt ở hai đầu đối diện của một đường kính và cho chúng chuyển động theo cùng một chiều quay. Các động lượng thông thường của chúng triệt tiêu nhau, nên khối tâm đứng yên, nhưng các mômen động lượng của chúng hướng theo cùng một chiều và cộng lại. Hệ có tổng động lượng bằng không và tổng mômen động lượng khác không: một mô hình cổ điển của spin.

Gọi vị trí của một hạt so với gốc là
\begin{equation}
\vec r=(x,y,z)=(x_1,x_2,x_3),
\end{equation}
và động lượng của nó là
\begin{equation}
\vec p=(p_x,p_y,p_z)=(p_1,p_2,p_3).
\end{equation}
Mômen động lượng là tích vectơ
\begin{equation}
\vec L=\vec r\times\vec p.
\end{equation}
Đảo thứ tự các thừa số thì vectơ đổi dấu: \(\vec p\times\vec r=-\vec r\times\vec p\). Quy ước chuẩn là \(\vec r\times\vec p\), với hướng của nó được xác định bởi quy tắc bàn tay phải.\footnote{\raggedright Khung hình video bài giảng \nolinkurl{lecture_07_figure_02.png}, 00:18:49.\par}

\begin{figure}[h]
\centering
\includegraphics[width=0.72\textwidth]{lecture_07_figure_02.png}
\par
\begin{tikzpicture}[scale=1.0]
\draw[->] (-0.6,0) -- (4.8,0) node[right] {$x$};
\draw[->] (0,-0.6) -- (0,3.8) node[above] {$y$};
\coordinate (P) at (3.2,2.3);
\draw[->,thick] (0,0) -- (P) node[midway,above left] {$\vec r$};
\draw[dashed] (P) -- (3.2,0);
\draw[dashed] (P) -- (0,2.3);
\draw[->,thick] (P) -- ++(-1.0,1.25) node[above] {$\vec p$};
\node at (4.15,2.75) {$\vec L=\vec r\times\vec p$};
\end{tikzpicture}
\caption{Đối với một hạt ở vị trí \(\vec r\) với động lượng \(\vec p\), mômen động lượng \(\vec L=\vec r\times\vec p\) vuông góc với mặt phẳng của chúng.}
\lectureframe{00:18:49}
\end{figure}

Trong các thành phần Đề-các,
\begin{equation}
\begin{aligned}
L_x&=yp_z-zp_y,\qquad
L_y=zp_x-xp_z,\\
L_z&=xp_y-yp_x.
\end{aligned}
\end{equation}
Chỉ cần nhớ công thức thứ nhất rồi chu chuyển các nhãn
\[
x\longrightarrow y\longrightarrow z\longrightarrow x.
\]
Chu chuyển cả nhãn thành phần lẫn các nhãn ở vế phải sẽ cho công thức tiếp theo mà không làm thay đổi dấu tương đối. Mẹo nhớ này hữu ích chính vì rất dễ vô tình đảo sai dấu.

Đối với một hạt điểm, \(\vec r\times\vec p\) là mômen động lượng quỹ đạo đối với gốc tọa độ đã chọn. Đối với một hệ hợp phần, hãy cộng các mômen động lượng của mọi thành phần theo quy tắc vectơ:
\begin{equation}
\vec L_{\mathrm{tot}}=\sum_a \vec r_a\times\vec p_a.
\end{equation}
Khi tổng này được tính trong hệ quy chiếu khối tâm, nó vẫn có thể khác không mặc dù
\[
\sum_a\vec p_a=0.
\]
Mômen động lượng còn lại đó là spin của hệ hợp phần.

\section{Toán tử lượng tử và các giao hoán tử chính tắc}

Trong cơ học lượng tử, các đại lượng quan sát được được biểu diễn bằng các toán tử. Vì vậy vị trí, động lượng và từng thành phần của mômen động lượng đều là các toán tử. Ở đây không cần đến cách hiện thực chi tiết của chúng dưới dạng toán tử vi phân. Các hệ thức giao hoán chứa đủ thông tin để suy ra phổ mômen động lượng.

Đối với hai toán tử \(A\) và \(B\),
\begin{equation}
[A,B]=AB-BA.
\end{equation}
Nếu hai đại lượng quan sát được giao hoán, chúng có thể được xác định đồng thời. Nếu chúng không giao hoán, cơ học lượng tử ngăn cản việc xác định chính xác đồng thời của chúng.

Cả ba tọa độ vị trí đều có thể được đo cùng nhau, nên các thành phần vị trí giao hoán:
\begin{equation}
[x_i,x_j]=0.
\end{equation}
Đặc biệt, mọi toán tử đều giao hoán với chính nó. Sẽ vô nghĩa nếu một đại lượng quan sát được lại không tương thích với một phép đo khác của chính đại lượng ấy.

Các thành phần động lượng cũng giao hoán:
\begin{equation}
[p_i,p_j]=0.
\end{equation}
Một thành phần vị trí giao hoán với một thành phần động lượng theo phương khác, nhưng không giao hoán với thành phần động lượng cùng phương với nó. Ba mệnh đề chính tắc có thể được viết gộp thành
\begin{equation}
[x_i,x_j]=0,\qquad
[p_i,p_j]=0,\qquad
[x_i,p_j]=i\hbar\,\delta_{ij}.
\end{equation}
Vì vậy, chẳng hạn,
\[
[x,p_y]=0,\qquad [x,p_x]=i\hbar,
\]
và đảo thứ tự thì đổi dấu:
\begin{equation}
[p_j,x_i]=-i\hbar\,\delta_{ij}.
\end{equation}

Mômen động lượng và hằng số Planck có cùng thứ nguyên. Mômen động lượng là khoảng cách nhân động lượng, còn hệ thức bất định có dạng về thứ nguyên
\begin{equation}
\Delta x\,\Delta p\sim\hbar
\end{equation}
ngoại trừ hệ số số học quen thuộc trong bất đẳng thức chính xác. Vì vậy, việc mômen động lượng sẽ bị lượng tử hóa theo đơn vị \(\hbar\) là điều tự nhiên, mặc dù sự lượng tử hóa ấy vẫn chưa được suy ra.

\section{Suy ra đại số mômen động lượng}

Nhiệm vụ đầu tiên là tính giao hoán tử của hai thành phần mômen động lượng. Từ
\[
L_x=yp_z-zp_y,\qquad
L_y=zp_x-xp_z,
\]
ta có
\begin{align}
[L_x,L_y]
&=[yp_z-zp_y,\;zp_x-xp_z] \\
&=[yp_z,zp_x]-[yp_z,xp_z] \notag\\
&\quad{}-[zp_y,zp_x]+[zp_y,xp_z].
\end{align}

\begin{classroomqa}
\lecturetimestamp{00:31:41}
\audiencequestion{Có thể nhớ các công thức thành phần bằng cách chu chuyển các ký hiệu từ trái sang phải thay vì chu chuyển ba phương trình từ trên xuống dưới không?}
\lecturerresponse{Có. Dùng mẹo nhớ kiểu chu chuyển nào cũng được. Mục đích duy nhất của mẹo nhớ là tạo ra các thành phần theo đúng thứ tự chu kỳ đồng thời theo dõi dấu.}
\end{classroomqa}

Có thể hiểu một giao hoán tử bằng cách hỏi xem các thừa số trong một tích có thể đổi chỗ qua các thừa số trong tích kia hay không. Trong
\[
[yp_z,xp_z],
\]
toán tử \(y\) giao hoán với cả \(x\) lẫn \(p_z\), còn \(p_z\) giao hoán với cả \(x\) lẫn chính nó. Không có sự cản trở nào, nên hạng này triệt tiêu. Lập luận tương tự cho
\[
[zp_y,zp_x]=0.
\]
Chỉ còn lại hai cặp tích.

Đối với cặp thứ nhất, trở ngại duy nhất là \(p_z\) đi qua \(z\):
\begin{align}
[yp_z,zp_x]
&=yp_zzp_x-zp_xyp_z \\
&=y(p_zz-zp_z)p_x \\
&=y[p_z,z]p_x \\
&=-i\hbar\,yp_x.
\end{align}
Thứ tự của giao hoán tử chính tắc là điều quan trọng:
\[
[p_z,z]=-[z,p_z]=-i\hbar.
\]

Đối với cặp còn lại không triệt tiêu, các dấu trừ trong định nghĩa của \(L_x\) và \(L_y\) nhân với nhau thành dấu cộng. Trở ngại duy nhất là \(z\) đi qua \(p_z\):
\begin{align}
[zp_y,xp_z]
&=zp_yxp_z-xp_zzp_y \\
&=x(zp_z-p_zz)p_y \\
&=x[z,p_z]p_y \\
&=i\hbar\,xp_y.
\end{align}
Gộp hai đóng góp lại,
\begin{align}
[L_x,L_y]
&=-i\hbar\,yp_x+i\hbar\,xp_y \\
&=i\hbar(xp_y-yp_x) \\
&=i\hbar L_z.
\end{align}

Đây là phép đơn giản hóa quan trọng. Lấy giao hoán tử của hai thành phần mômen động lượng sẽ tạo ra thành phần thứ ba; không còn toán tử vị trí hay động lượng nào nữa. Hoán vị theo chu kỳ cho
\begin{equation}
\begin{aligned}
[L_x,L_y]&=i\hbar L_z,\qquad
[L_y,L_z]=i\hbar L_x,\\
[L_z,L_x]&=i\hbar L_y.
\end{aligned}
\end{equation}
Tương đương,
\begin{equation}
[L_i,L_j]=i\hbar\,\epsilon_{ijk}L_k.
\end{equation}

Đại số này khép kín trên chính các toán tử mômen động lượng. Một khi đã thu được các hệ thức này, \(x_i\) và \(p_i\) dùng để suy ra chúng có thể được gác sang một bên. Cùng đại số ấy áp dụng khi tổng mômen động lượng là một tổng trên nhiều thành phần cấu thành. Đặc biệt, nó áp dụng trong hệ quy chiếu khối tâm, nơi mômen động lượng còn lại là spin. Vì thế các giao hoán tử mômen động lượng có thể được dùng một cách trừu tượng mà không cần trước hết dựng spin như chuyển động của những phần cấu thành nhỏ.

Việc chọn các trục \(x\), \(y\) và \(z\) là tùy ý. Hãy quay sang bất kỳ bộ trục Đề-các quán tính trực giao nào khác, và các thành phần có dấu phẩy sẽ tuân theo cùng những hệ thức:
\begin{equation}
[L'_x,L'_y]=i\hbar L'_z,
\end{equation}
cùng với các hệ thức thu được bằng hoán vị theo chu kỳ. Đại số là bất biến quay, nên các dự đoán vật lý không thể phụ thuộc vào việc đã chọn những trục quán tính nào cho phép tính.

\begin{classroomqa}
\lecturetimestamp{00:39:09}
\audiencequestion{Mômen động lượng spin có độc lập với các phép tịnh tiến cũng như các phép quay của các trục tọa độ không?}
\lecturerresponse{Có. Spin nội tại độc lập với phép tịnh tiến. Mômen động lượng quỹ đạo đối với một gốc có thể thay đổi khi gốc đó được tịnh tiến, nhưng spin còn lại trong hệ quy chiếu khối tâm thì không phụ thuộc vào phép tịnh tiến ấy.}
\end{classroomqa}

Đại số trừu tượng có thể sinh ra vô số đồng nhất thức, phần lớn trong số đó không đáng quan tâm về mặt vật lý. Giá trị của nó bộc lộ khi một đồng nhất thức trở thành một sự kiện có thể đo được hoặc hé lộ điều gì đó về ý nghĩa thao tác của phép đo. Cho đến thời điểm này, phép xây dựng vẫn chủ yếu mang tính đại số. Bước tiếp theo rút ra cấu trúc rời rạc của spin.

\begin{classroomqa}
\lecturetimestamp{00:40:42}
\audiencequestion{Nếu các trục \(x\), \(y\) và \(z\) đã chọn quay cùng các hạt thì sao?}
\lecturerresponse{Phần thảo luận này chỉ giới hạn trong các hệ quy chiếu quán tính. Một hệ tọa độ quay là phi quán tính và đưa vào những hiệu ứng như lực ly tâm và lực Coriolis. Mômen động lượng không bất biến khi chuyển sang một hệ quy chiếu cùng quay: một hệ đang quay thấy rõ trong một hệ quán tính có thể trông như đứng yên trong một hệ quy chiếu quay cùng nó. Tính bất biến quay được dùng ở đây liên quan đến các trục Cartesian quán tính có định hướng khác nhau, chứ không phải các trục quay theo thời gian.}
\end{classroomqa}

\section{Toán tử bậc thang và trục lượng tử hóa}

Đưa vào hai tổ hợp
\begin{equation}
L_+=L_x+iL_y,
\qquad
L_-=L_x-iL_y.
\end{equation}
Chúng tương tự các toán tử tạo và hủy của dao động tử điều hòa. Vai trò của chúng là nâng và hạ một thành phần được chọn của mômen động lượng.

Không có gì đặc biệt về trục \(z\), nhưng phải chọn một trục nếu muốn chỉ ra một thành phần. Lấy \(z\) làm trục lượng tử hóa và viết
\begin{equation}
L_z\lvert m\rangle=m\hbar\,\lvert m\rangle.
\end{equation}
Do đó \(m\) là thành phần \(z\) được đo theo đơn vị \(\hbar\). Trong khi thực hiện đại số, người ta thường đặt \(\hbar=1\); ở đây nó sẽ được giữ tường minh.

Giao hoán tử \([L_+,L_-]\) tỉ lệ với \(L_z\), nhưng nó không cần cho lập luận trước mắt. Điều trước hết quan trọng là giao hoán tử của mỗi toán tử bậc thang với \(L_z\). Từ
\[
[L_x,L_z]=-i\hbar L_y,
\qquad
[L_y,L_z]=i\hbar L_x,
\]
ta thu được
\begin{align}
[L_+,L_z]
&=[L_x+iL_y,L_z] \\
&=[L_x,L_z]+i[L_y,L_z] \\
&=-i\hbar L_y+i(i\hbar L_x) \\
&=-\hbar(L_x+iL_y) \\
&=-\hbar L_+.
\end{align}
Đổi dấu của hạng \(iL_y\) thì được
\begin{equation}
[L_-,L_z]=+\hbar L_-.
\end{equation}
Tương đương, nếu đảo thứ tự của mỗi giao hoán tử,
\begin{equation}
[L_z,L_+]=\hbar L_+,
\qquad
[L_z,L_-]=-\hbar L_-.
\end{equation}

Tại một thời điểm chỉ một thành phần của mômen động lượng có thể có giá trị xác định chính xác. Bất kỳ hai thành phần phân biệt nào cũng không giao hoán, nên chúng không thể được xác định đồng thời. Thành phần xác định được chọn định nghĩa trục lượng tử hóa; việc chọn \(z\) chỉ là thuận tiện, không phải một ưu tiên vật lý.

Bây giờ hãy tác dụng \(L_+\) lên một trạng thái riêng của \(L_z\). Có thể viết lại giao hoán tử thành
\begin{equation}
L_zL_+=L_+L_z+\hbar L_+.
\end{equation}
Do đó
\begin{align}
L_z\bigl(L_+\lvert m\rangle\bigr)
&=\bigl(L_+L_z+\hbar L_+\bigr)\lvert m\rangle \\
&=m\hbar L_+\lvert m\rangle+\hbar L_+\lvert m\rangle \\
&=(m+1)\hbar\,L_+\lvert m\rangle.
\end{align}
Miễn là nó khác không, \(L_+\lvert m\rangle\) là một trạng thái riêng của \(L_z\) với trị riêng đã tăng thêm một đơn vị \(\hbar\). Tương tự,
\begin{equation}
L_z\bigl(L_-\lvert m\rangle\bigr)
=(m-1)\hbar\,L_-\lvert m\rangle.
\end{equation}
Vì vậy \(L_-\) hạ \(m\) xuống một đơn vị.

Đây đã là một kết luận không tầm thường. Dù phổ của \(L_z\) rốt cuộc là gì, các giá trị lân cận sai khác nhau bởi các số nguyên:
\begin{equation}
\ldots,\ m-2,\ m-1,\ m,\ m+1,\ m+2,\ \ldots
\end{equation}
Các đầu mút vẫn chưa được xác định.

\section{Các đa bội hữu hạn và nhãn spin}

Một loại hạt cố định có độ lớn mômen động lượng cố định. Mômen động lượng của nó có thể được định hướng lại, nhưng nếu tăng độ lớn của nó mà không có giới hạn thì cuối cùng sẽ tạo ra một đối tượng khác. Vì thế phải có một thành phần \(z\) cực đại khả dĩ đối với hạt đó.

Làm sao một bậc thang có toán tử nâng lại có đỉnh? Dao động tử điều hòa cung cấp mô hình. Toán tử hủy của nó cho kết quả bằng không khi chạm đến trạng thái đáy. Ở đây toán tử nâng phải cho bằng không khi chạm đến trạng thái đỉnh:
\begin{equation}
L_+\lvert m_{\max}\rangle=0.
\end{equation}
Phát biểu rằng \(L_+\) làm tăng một trị riêng chỉ áp dụng khi vectơ nó tạo ra là khác không. Trạng thái trọng số cao nhất chính là trạng thái ngoại lệ mà tại đó dãy nâng chấm dứt.

Tính bất biến quay cố định đầu mút dưới. Nếu mômen động lượng có thể hướng cực đại dọc theo \(+z\), thì cùng đối tượng vật lý ấy có thể được quay sao cho nó hướng cực đại dọc theo \(-z\). Do đó
\begin{equation}
m_{\min}=-m_{\max}.
\end{equation}
Phổ đối xứng qua 0 và có khoảng cách đơn vị. Có hai họ khả dĩ.

Trong họ thứ nhất, có giá trị 0:
\begin{equation}
\begin{aligned}
m={}&-m_{\max},-m_{\max}+1,\ldots,-1,0,\\
&1,\ldots,m_{\max}-1,m_{\max}.
\end{aligned}
\end{equation}
Khi đó \(m_{\max}\) là một số nguyên, và số trạng thái là
\begin{equation}
2m_{\max}+1.
\end{equation}
Giá trị cực đại đặc trưng cho hạt và được gọi là spin của nó. Theo ký hiệu quy ước,
\begin{equation}
s\equiv m_{\max},
\qquad
\lvert s,m\rangle\equiv\lvert m\rangle
\end{equation}
trong một đa bội có spin cố định.\footnote{\raggedright Khung hình video bài giảng \nolinkurl{lecture_07_figure_03.png}, 01:01:08.\par}

\begin{figure}[htbp]
\centering
\includegraphics[width=0.70\textwidth]{lecture_07_figure_03.png}
\caption{Điều kiện trọng số cao nhất chấm dứt dãy nâng, và các giá trị \(m\) được phép tạo thành một bậc thang hữu hạn với khoảng cách đơn vị.}
\lectureframe{01:01:08}
\end{figure}

Một hạt spin-0 chỉ có \(m=0\). Một hạt spin-1 có
\[
m=1,0,-1,
\]
và một hạt spin-2 có 5 giá trị thành phần. Boson Higgs là một ví dụ về spin 0. Photon và boson \(Z\) có spin 1. Graviton, nếu tồn tại, có spin 2. Các đối tượng hợp thành có thể có spin 3 và lớn hơn nhiều: hạt nhân có thể như vậy, và những vật quay đủ vĩ mô như quả bóng rổ và các thiên hà có thể mang mômen động lượng lớn tùy ý. Điều vắng mặt trong danh sách các hạt sơ cấp thông thường là một hạt sơ cấp spin-3.

Họ đối xứng còn lại, cũng có khoảng cách đơn vị, không chứa 0. Nó bắt đầu bằng
\[
m=\frac12,-\frac12
\]
và kéo dài theo các bước đơn vị:
\begin{equation}
\begin{aligned}
m={}&-m_{\max},-m_{\max}+1,\ldots,-\frac32,-\frac12,\\
&\frac12,\frac32,\ldots,m_{\max}.
\end{aligned}
\end{equation}
Ở đây \(m_{\max}\) là một bán số nguyên. Trường hợp đơn giản nhất có spin \(\tfrac12\), tiếp theo là spin \(\tfrac32\), spin \(\tfrac52\), v.v. Vì vậy các đa bội mômen động lượng thuộc hai họ: spin nguyên và spin bán nguyên.

\begin{classroomqa}
\lecturetimestamp{01:03:46}
\audiencequestion{Một đa bội spin bán nguyên cũng có một giá trị cực đại chứ?}
\lecturerresponse{Có. Bậc thang một lần nữa lại có một giá trị cao nhất. Đối với hạt spin-\(\tfrac32\), giá trị cực đại là \(m_{\max}=\tfrac32\); đối với hạt spin-\(\tfrac52\), nó là \(m_{\max}=\tfrac52\). Khác biệt duy nhất là giá trị cực đại và mọi giá trị khác trong đa bội đều là bán số nguyên.}
\end{classroomqa}

\begin{classroomqa}
\lecturetimestamp{01:04:04}
\audiencequestion{Có phải không có hạt sơ cấp nào có spin 3 hoặc cao hơn?}
\lecturerresponse{Trong khuôn khổ lượng tử tương đối tính đang xét, các spin hạt sơ cấp trong danh sách thông thường là \(0,\tfrac12,1,\tfrac32,\) và \(2\). Các đối tượng hợp thành không bị giới hạn trong khoảng này và có thể có spin 3 hoặc spin lớn hơn nhiều.}
\end{classroomqa}

\begin{classroomqa}
\lecturetimestamp{01:04:50}
\audiencequestion{Các giá trị âm trên bậc thang có biểu thị spin âm không?}
\lecturerresponse{Không. Một \(m\) âm là một giá trị thành phần: nó cho biết mômen động lượng hướng ngược lại so với trục đã chọn. Spin của hạt là giá trị cực đại không âm \(s=m_{\max}\).}
\end{classroomqa}

\begin{classroomqa}
\lecturetimestamp{01:05:10}
\audiencequestion{Nếu các giá trị thành phần là rời rạc, vậy mọi định hướng ở giữa chúng thì sao?}
\lecturerresponse{Không phải chỉ có 2 hay 3 định hướng không gian khả dĩ. Các định hướng trung gian được biểu diễn bởi các chồng chập lượng tử của các trạng thái cơ sở \(m\). Cấu trúc spin-\(\tfrac12\) bên dưới sẽ làm điều này tường minh.}
\end{classroomqa}

\begin{classroomqa}
\lecturetimestamp{01:05:42}
\audiencequestion{Còn chuyển động kiểu vít xoắn hoặc sóng phân cực tròn thì sao?}
\lecturerresponse{Một sóng điện từ phân cực tròn được cấu tạo từ các photon phân cực tròn, và các photon đó mang mômen động lượng dọc theo trục truyền. Từ ``vít xoắn'' kết hợp sự quay với chuyển động về phía trước, nên đó không phải là mô tả đúng cho một đối tượng đang được xét ở trạng thái nghỉ. Đối với một hạt đang chuyển động, đại lượng liên quan là độ xoắn: định hướng spin của nó so với hướng chuyển động của nó.}
\end{classroomqa}

\begin{classroomqa}
\lecturetimestamp{01:07:10}
\audiencequestion{Có hạt sơ cấp đã biết nào có spin \(\tfrac32\) không?}
\lecturerresponse{Chưa có hạt sơ cấp spin-\(\tfrac32\) nào được xác nhận trực tiếp. Siêu đối xứng giả thuyết một đối tác spin-\(\tfrac32\) của graviton gọi là gravitino, nhưng nó chưa được phát hiện trực tiếp.}
\end{classroomqa}

\begin{classroomqa}
\lecturetimestamp{01:08:01}
\audiencequestion{Có phải các toán tử mômen động lượng đã thôi không còn tác dụng lên không gian hàm sóng hay không gian ket ban đầu nữa không?}
\lecturerresponse{Không. Chưa có không gian spin ban đầu nào được xác định cả. Đại số đang được dùng để xác định không gian đó. Với spin cố định \(s\), nó có các vectơ cơ sở được gắn nhãn bởi các giá trị cho phép
\[
\lvert s,m\rangle,\qquad m=-s,-s+1,\ldots,s.
\]
Không gian spin nội tại này khác với không gian thông thường của các hàm sóng không gian, nơi mômen động lượng quỹ đạo tác dụng như \(\mathbf r\times\mathbf p\), mặc dù cả hai đều hiện thực cùng một đại số mômen động lượng.}
\end{classroomqa}

\section{Độ lớn của mômen động lượng}

\begin{classroomqa}
\lecturetimestamp{01:09:32}
\audiencequestion{Độ lớn của vectơ mômen động lượng có phải là đại lượng quan sát được không?}
\lecturerresponse{Có. Thuận tiện nhất là dùng bình phương của nó,
\[
L^2=L_x^2+L_y^2+L_z^2.
\]
Giá trị của nó trên trạng thái trọng số cao nhất không đơn giản chỉ là \(s^2\hbar^2\). Ngay cả khi \(L_z\) có giá trị xác định và cực đại, các thành phần ngang không giao hoán vẫn không thể đồng thời xác định chính xác bằng 0, nên chúng cung cấp một đóng góp lượng tử bổ sung.}
\end{classroomqa}

Để tính đóng góp đó, hãy nhân các toán tử bậc thang theo thứ tự sao cho toán tử nâng tác dụng trước lên một ket:
\begin{align}
L_-L_+
&=(L_x-iL_y)(L_x+iL_y) \\
&=L_x^2+L_y^2+iL_xL_y-iL_yL_x \\
&=L_x^2+L_y^2+i[L_x,L_y] \\
&=L_x^2+L_y^2-\hbar L_z.
\end{align}
Bước cuối cùng dùng \([L_x,L_y]=i\hbar L_z\). Do đó
\begin{equation}
L_x^2+L_y^2=L_-L_++\hbar L_z,
\end{equation}
và vì vậy
\begin{equation}
L^2=L_-L_++L_z^2+\hbar L_z.
\end{equation}

Áp dụng đồng nhất thức này lên trạng thái trọng số cao nhất \(\lvert s,s\rangle=\lvert m_{\max}\rangle\):
\begin{align}
L^2\lvert s,s\rangle
&=\left(L_-L_++L_z^2+\hbar L_z\right)\lvert s,s\rangle \\
&=L_-L_+\lvert s,s\rangle
+s^2\hbar^2\lvert s,s\rangle
+s\hbar^2\lvert s,s\rangle.
\end{align}
Hạng đầu tiên triệt tiêu vì \(L_+\lvert s,s\rangle=0\). Do đó
\begin{equation}
L^2\lvert s,s\rangle
=s(s+1)\hbar^2\lvert s,s\rangle.
\end{equation}
Phần \(+s\) bổ sung là hiệu chỉnh lượng tử gắn với những thăng giáng ngang không thể tránh khỏi. Theo cổ điển, ta chỉ mong đợi \(s^2\hbar^2\).

Một vài giá trị đầu tiên làm cho sự khác biệt trở nên cụ thể:
\begin{align}
s=0 &: \qquad L^2=0,\\
s=1 &: \qquad L^2=2\hbar^2,\\
s=\frac12 &: \qquad L^2=\frac34\hbar^2.
\end{align}
Với \(s=1000\), hệ số không thứ nguyên là \(1000\times1001\), khác rất ít so với \(1000^2\) trên thang của một mômen động lượng lớn như vậy. Nói chung,
\begin{equation}
s(s+1)=s^2\left(1+\frac1s\right),
\end{equation}
nên hiệu chỉnh lượng tử tương đối trở nên không đáng kể đối với \(s\) rất lớn. Đó là cách kết quả cổ điển \(s^2\) xuất hiện.

Mọi trạng thái \(m\) trong một đa bội spin cố định đều có cùng trị riêng này của \(L^2\). Thay đổi \(m\) làm đổi thành phần dọc theo trục đã chọn, chứ không làm đổi độ lớn toàn phần. Đây là phiên bản lượng tử của việc quay một vectơ mômen động lượng mà không thay đổi độ dài của nó.

\begin{classroomqa}
\lecturetimestamp{01:16:43}
\audiencequestion{Có thể biết đồng thời \(L^2\) và thành phần \(z\) không?}
\lecturerresponse{Có. Bình phương toàn phần giao hoán với mọi thành phần của mômen động lượng:
\[
[L^2,L_i]=0,\qquad i=x,y,z.
\]
Vì vậy \(L^2\) và một thành phần được chọn, chẳng hạn \(L_z\), có thể được xác định cùng nhau mặc dù hai thành phần phân biệt thì không thể. Chẳng hạn, mọi electron đều có \(L^2=\tfrac34\hbar^2\), bất kể định hướng của nó.}
\end{classroomqa}

\begin{classroomqa}
\lecturetimestamp{01:17:41}
\audiencequestion{Có phải \(L^2\) bằng 0 đối với spin 1, hay chỉ đối với spin 0?}
\lecturerresponse{Chỉ spin 0 mới có \(L^2=0\). Spin 1 cho
\[
L^2=1(1+1)\hbar^2=2\hbar^2,
\]
trong khi spin 3 cho
\[
L^2=3(3+1)\hbar^2=12\hbar^2.
\]}
\end{classroomqa}

\begin{classroomqa}
\lecturetimestamp{01:18:13}
\audiencequestion{Vì sao mọi mức \(m\) trong đa bội đều có cùng trị riêng \(L^2\)?}
\lecturerresponse{Hãy bắt đầu từ trạng thái trọng số cao nhất rồi hạ nó xuống. Vì \(L^2\) giao hoán với mọi thành phần, nó cũng giao hoán với \(L_-\):
\[
[L^2,L_-]=0.
\]
Nếu
\[
L^2\lvert s,s\rangle=s(s+1)\hbar^2\lvert s,s\rangle,
\]
thì
\begin{align}
L^2\bigl(L_-\lvert s,s\rangle\bigr)
&=L_-L^2\lvert s,s\rangle \\
&=s(s+1)\hbar^2\,L_-\lvert s,s\rangle.
\end{align}
Vì vậy trạng thái ngay bên dưới có cùng trị riêng \(L^2\). Lặp lại lập luận cho thấy trị riêng đó giữ nguyên suốt bậc thang.}
\end{classroomqa}

Vì vậy, các giá trị \(m\) rời rạc không có nghĩa là chỉ tồn tại một tập rời rạc các định hướng không gian. Tính bất biến quay vẫn đòi hỏi các trạng thái tương ứng với những hướng thu được bằng phép quay. Đối với spin \(\tfrac12\), những hướng ấy sẽ xuất hiện dưới dạng các chồng chập của hai trạng thái cơ sở \(L_z\).

\section{Trạng thái spin-\texorpdfstring{\(\tfrac12\)}{1/2} và định hướng}

Tính bất biến quay đòi hỏi nhiều hơn là các trạng thái riêng rời rạc của \(L_z\). Một spin có thể hướng theo bất kỳ trục nào, dù phép đo thành phần của nó dọc theo một trục đã chọn chỉ cho các giá trị rời rạc được phép. Các hướng nằm giữa các trạng thái riêng theo trục \(z\) được biểu diễn bởi các chồng chập.

Đối với một hạt spin-\(\tfrac12\), chỉ xét các bậc tự do spin. Electron còn có các trạng thái quỹ đạo, nên chắc chắn nó không phải tổng cộng chỉ có hai trạng thái; riêng không gian spin của nó là hai chiều. Chọn cơ sở theo trục \(z\)
\begin{equation}
L_z\lvert+\rangle=\frac{\hbar}{2}\lvert+\rangle,
\qquad
L_z\lvert-\rangle=-\frac{\hbar}{2}\lvert-\rangle.
\end{equation}
Một trạng thái spin tổng quát là
\begin{equation}
\lvert\psi\rangle
=
\alpha\lvert+\rangle+\beta\lvert-\rangle,
\end{equation}
trong đó \(\alpha\) và \(\beta\) là các số phức. Phép đo \(L_z\) cho \(+\hbar/2\) với xác suất \(|\alpha|^2\) và \(-\hbar/2\) với xác suất \(|\beta|^2\), nên điều kiện chuẩn hóa đòi hỏi
\begin{equation}
|\alpha|^2+|\beta|^2=1.
\end{equation}

Một pha chung không có tác dụng vật lý:
\begin{equation}
(\alpha,\beta)\longrightarrow
e^{i\chi}(\alpha,\beta).
\end{equation}
Hai biên độ phức ban đầu chứa bốn tham số thực. Điều kiện chuẩn hóa loại bỏ một tham số, và pha toàn cục không liên quan loại bỏ thêm một tham số nữa, nên còn lại hai tham số thực vật lý. Đó chính là số tham số cần để xác định một hướng bằng góc cực và góc phương vị của nó. Khi thay đổi \(\alpha\) và \(\beta\), ta làm quay hướng spin.

Một thành phần được đo và giá trị kỳ vọng của nó là hai thứ khác nhau. Giá trị kỳ vọng là trung bình trên một tập hợp. Hãy chuẩn bị nhiều electron ở đúng cùng một trạng thái, đo \(L_z\) một lần trên mỗi electron, cộng các kết quả lại, rồi chia cho số thí nghiệm. Mỗi kết quả riêng lẻ vẫn là \(+\hbar/2\) hoặc \(-\hbar/2\), nhưng giá trị trung bình có thể nhận bất kỳ giá trị nào giữa chúng:
\begin{equation}
\langle L_z\rangle
=
\frac{\hbar}{2}\bigl(|\alpha|^2-|\beta|^2\bigr).
\end{equation}

Các biên độ hướng lên và hướng xuống bằng nhau cho một spin định hướng theo một trong hai chiều của trục \(x\). Các đại diện đã chuẩn hóa là\footnote{Tỉ số giữa các hệ số xác định các hướng. Các thừa số \(1/\sqrt2\) cung cấp sự chuẩn hóa do \( |\alpha|^2+|\beta|^2=1\) đòi hỏi. Các pha tương đối \(+i\) và \(-i\) cho hai hướng đối nhau dọc theo trục \(y\); việc gán dấu nào cho hướng \(y\) dương phụ thuộc vào quy ước trục.}
\begin{equation}
\lvert x;+\rangle
=
\frac{\lvert+\rangle+\lvert-\rangle}{\sqrt2},
\qquad
\lvert x;-\rangle
=
\frac{\lvert+\rangle-\lvert-\rangle}{\sqrt2}.
\end{equation}
Đối với một trong hai trạng thái, phép đo theo \(z\) cho hai kết quả khả dĩ với xác suất bằng nhau, và vì thế
\begin{equation}
\langle L_z\rangle
=
\frac12\left(\frac{\hbar}{2}\right)
+
\frac12\left(-\frac{\hbar}{2}\right)
=0.
\end{equation}
Dấu tương đối cộng và trừ cho hai hướng đối nhau dọc theo trục \(x\). Tương tự,
\begin{equation}
\frac{\lvert+\rangle+i\lvert-\rangle}{\sqrt2},
\qquad
\frac{\lvert+\rangle-i\lvert-\rangle}{\sqrt2}
\end{equation}
cho hai hướng đối nhau của trục \(y\). Các giá trị chuẩn hóa tổng quát của \(\alpha\) và \(\beta\) mô tả các trục tùy ý. Nói rằng một spin hướng dọc theo một trục cụ thể có nghĩa là phép đo thành phần theo trục đó cho một kết quả xác định. So với một trục khác, cùng trạng thái ấy nói chung là một chồng chập.

Các kết quả đo riêng lẻ vẫn được lượng tử hóa. Chính các giá trị trung bình trên tập hợp của chúng mới biến thiên liên tục và tái hiện hình học quen thuộc của một vectơ cổ điển.

\section{Spin, thống kê và các hạt hợp thành}

Khối lượng, điện tích và spin là một số trong những tính chất chính dùng để đặc trưng một hạt. Spin còn tương quan với thống kê. Quan hệ spin-thống kê là một định lý của cơ học lượng tử tương đối tính, mặc dù ở đây sẽ không chứng minh:
\begin{equation}
\begin{aligned}
\text{half-integer spin} &\quad\Longleftrightarrow\quad \text{fermion},\\
\text{integer spin} &\quad\Longleftrightarrow\quad \text{boson}.
\end{aligned}
\end{equation}
Hai fermion đồng nhất không thể chiếm cùng một trạng thái lượng tử hoàn chỉnh. Các boson đồng nhất có thể chiếm cùng một trạng thái, và số lượng lớn boson trong một trạng thái có thể cùng nhau hành xử như một sóng cổ điển. Vì vậy, các hạt spin \(\tfrac12,\tfrac32,\tfrac52,\ldots\) là fermion, còn các hạt spin \(0,1,2,\ldots\) là boson.

Quy tắc này mở rộng ngay sang các hệ hợp thành. Cộng hai mômen góc bán nguyên cho ra một mômen góc nguyên. Tổng quát hơn, một số chẵn thành phần fermion tạo thành một hạt hợp thành boson, còn một số lẻ tạo thành một hạt hợp thành fermion. Việc thêm các thành phần boson không làm thay đổi quy tắc chẵn lẻ này.

Hydro thông thường chứa một proton spin-\(\tfrac12\) và một electron spin-\(\tfrac12\). Vì vậy, xét tổng thể nó chứa một số chẵn fermion và là boson. Một deuteron chứa một proton và một neutron, cả hai đều là fermion, nên cũng là boson. Một nguyên tử deuteri thêm một electron vào deuteron, cho ba thành phần fermion ở mức mô tả này, nên deuteri là fermion.

Hạt và phản hạt có cùng khối lượng và spin, nhưng điện tích trái dấu. Positronium là một electron và một positron liên kết với nhau. Nó chứa hai fermion và vì thế là boson.

\begin{classroomqa}
\lecturetimestamp{01:33:08}
\audiencequestion{Các boson có đi xuyên qua vật chất không, và hiện tượng siêu dẫn có liên quan đến việc các cặp electron hành xử như boson không?}
\lecturerresponse{Tính chất liên quan không đơn thuần là khả năng xuyên qua vật chất; đó là khả năng của boson cùng chiếm một trạng thái. Heli siêu lỏng cho một ví dụ đơn giản hơn. Một nguyên tử heli-\(4\) thông thường chứa hai proton, hai neutron và hai electron, tức là một số chẵn thành phần fermion, nên nguyên tử ấy là boson. Nhiều nguyên tử heli có thể cùng chiếm một trạng thái lượng tử và chuyển động tập thể. Khi đó có thể hình dung gần đúng một chất siêu lỏng như một sóng cổ điển được hình thành bằng cách dồn nhiều nguyên tử vào cùng một trạng thái. Hình dung đó không hoàn toàn chính xác, nhưng nó nắm được ý chính cần cho chỗ này. Một electron trong kim loại là fermion và không thể chia sẻ một trạng thái hoàn chỉnh với một electron khác. Tuy nhiên, các cặp electron là các hợp thể boson, và đó là mối liên hệ với siêu dẫn. Điều này làm nảy sinh một nghịch lý: nếu các electron thành phần không thể cùng chiếm một trạng thái, thì làm sao nhiều cặp hợp thành lại có thể cùng chiếm một trạng thái? Cách giải quyết đòi hỏi phải trình bày cẩn thận hơn về các trạng thái hợp thành và sẽ được để lại sau.}
\end{classroomqa}

\begin{classroomqa}
\lecturetimestamp{01:35:46}
\audiencequestion{Vì proton và neutron là các hạt hợp thành, tại sao chúng lại có spin bán nguyên?}
\lecturerresponse{Mỗi hạt được tạo thành từ một số lẻ quark spin-\(\tfrac12\). Ba thành phần fermion tạo thành một hạt hợp thành fermion, nên proton và neutron có thể có spin bán nguyên.}
\end{classroomqa}

\section{Nguyên lý loại trừ Pauli và các obitan nguyên tử}

Bảng tuần hoàn cung cấp động lực trực tiếp cho nguyên lý loại trừ Pauli. Quy tắc không phải là hai electron không bao giờ có thể có cùng chuyển động quỹ đạo. Mà là hai electron không thể chiếm cùng một trạng thái lượng tử hoàn chỉnh.

Phân biệt này quan trọng ngay cả trong những nguyên tử đơn giản nhất. Bắt đầu từ hydro, tăng điện tích hạt nhân và thêm một electron thứ hai. Hai electron có thể chiếm cùng một obitan không gian, nhưng các trạng thái spin của chúng phải khác nhau. Một electron có thể spin lên và electron kia spin xuống so với trục đã chọn. Khi hai electron chiếm một obitan, spin của chúng kết hợp thành trạng thái singlet của spin, nên tổng spin của chúng triệt tiêu. Không quá hai electron có thể chiếm một obitan không gian đơn lẻ vì chỉ có hai trạng thái spin độc lập khả dụng.

\begin{classroomqa}
\lecturetimestamp{01:37:53}
\audiencequestion{Có thể hình dung hai spin đối nhau của hai electron trong một obitan như các bánh răng ăn khớp không?}
\lecturerresponse{Như một hình dung thô thì có: các bánh răng ăn khớp quay theo hai chiều ngược nhau. Đó chỉ là một phép loại suy cho các hướng spin đối ngược, không phải một mô hình cơ học theo nghĩa đen của các electron.}
\end{classroomqa}

Từ ``obitan'' có thể được dùng ở hai mức độ khác nhau. Một trạng thái chuyển động quỹ đạo, theo nghĩa rộng, là trạng thái thu được khi giải phương trình Schrödinger mà bỏ qua spin. Nó bao gồm nhiều thông tin hơn việc chỉ đơn thuần electron chiếm lớp nào.

\begin{classroomqa}
\lecturetimestamp{01:38:09}
\audiencequestion{Một obitan có nghĩa nhiều hơn việc chỉ đơn thuần electron chiếm lớp nào không?}
\lecturerresponse{Có. Ở đây, chuyển động quỹ đạo có nghĩa là trạng thái Schrödinger của electron khi bỏ qua spin. Nó bao gồm số lượng tử chính \(n\), số mang thông tin hướng tâm và xấp xỉ xác định khoảng cách của electron tới hạt nhân, cũng như các số lượng tử mômen góc quỹ đạo. Trong ký hiệu phổ học mà các nhà hóa học dùng, \(S\), \(P\), \(D\) và \(F\) lần lượt tương ứng với mômen góc quỹ đạo \(l=0,1,2,3\).}
\end{classroomqa}

Vì vậy, số lượng tử chính và mômen góc quỹ đạo mô tả những khía cạnh khác nhau của chuyển động. Các nhãn \(S,P,D,F\) nói tới giá trị của \(l\), còn \(n\) cung cấp thông tin hướng tâm riêng biệt. Quy tắc của Pauli là không có hai electron nào được phép có tất cả các số lượng tử giống hệt nhau, kể cả số lượng tử spin.

\begin{classroomqa}
\lecturetimestamp{01:40:42}
\audiencequestion{Các số lượng tử có thể xác định có bao nhiêu electron được phép chiếm mỗi lớp trong một nguyên tử không?}
\lecturerresponse{Có. Chúng xác định số tối đa có thể có, dù không xác định một nguyên tử cụ thể thực sự có bao nhiêu electron trong lớp đó. Bài toán đếm lớp tổng quát sẽ không được theo đuổi ở đây.}
\end{classroomqa}

Với một mômen góc quỹ đạo cố định \(l\), số lượng tử thành phần có
\begin{equation}
2l+1
\end{equation}
giá trị khả dĩ. Đây chính là phép đếm mômen góc đã được dùng cho spin. Một trạng thái \(S\) có \(l=0\), nên nó có một trạng thái obitan. Một trạng thái \(P\) có \(l=1\), nên nó có ba trạng thái obitan. Một trạng thái \(D\) có \(l=2\), nên nó có năm trạng thái obitan.

Spin làm gấp đôi mỗi số đếm obitan, vì mỗi obitan không gian có thể chứa một electron cho mỗi một trong hai trạng thái spin đối nhau.

\begin{classroomqa}
\lecturetimestamp{01:41:54}
\audiencequestion{Có phải trạng thái \(S\) chứa hai electron, và ba trạng thái \(P\) chứa thêm sáu electron nữa không?}
\lecturerresponse{Đúng. Trạng thái obitan \(S\) duy nhất chứa tối đa hai electron. Ba trạng thái obitan \(P\) mỗi trạng thái chứa hai electron, thành ra thêm sáu. Vì vậy, các trạng thái \(S\) và \(P\) đang bàn đến chứa tổng cộng \(2+6=8\) electron.}
\end{classroomqa}

Cùng cấu trúc mômen góc tổ chức các obitan nguyên tử ấy cũng xuất hiện trở lại, với một vài khác biệt, khi các quark chuyển động trong các hệ liên kết và quay quanh nhau.